% Pillar-7: data substrates and statistical conventions actually used by this
% paper. Replaces the shared TinkerRL-Bench methods block (task taxonomy,
% cross-library rosters, PPO hyperparameters, legacy BH table), none of which
% is analysed here (Deep Review 2026-10-18, findings 418--425).
\section{Data Substrates and Statistical Conventions}
\label{sec:p7-methods}
\label{sec:benchmark}
\label{sec:setup}
This paper re-analyses existing logs and reports one small live pilot. It
uses only the substrates below; the broader benchmark task suite, library
rosters, and hyperparameter tables of the companion papers are not used.
Every substrate has binary verifiable rewards $R\in\{0,1\}$.

\paragraph{S1: 368-run W\&B audit.} Project \texttt{zvf-audit} (368 runs; 350
with a run-level summary, 352 with per-step histories, 2{,}004 step rows),
nine model labels (the task is not recorded in the run configs). Batch size (6--32), $G$ (4--16) and step
budget (2--12) vary by model and are recorded per run
(\texttt{experiments/results/publication\_fixes/zvf\_confound/}). It is an
observational corpus, not a designed experiment.

\paragraph{S2: group-size sweep.} Qwen2.5-0.5B (base), LoRA ($r{=}16$,
$\alpha{=}32$, q/v projections), Adam at $10^{-4}$, two-operand addition with
operands drawn uniformly from $\{1,\dots,99\}$, 16 fresh prompts per step for
40 steps, sampling at temperature 1.0 with at most 10 new tokens; reward 1 iff
the last integer in the completion equals the sum. $G\in\{2,4,8,16\}\times$
seeds $\{42,123,456\}$ (12 runs); held-out accuracy is greedy on 200 fresh
problems (\texttt{experiments/modal/modal\_groupsize\_zvf\_sweep.py},
\texttt{experiments/results/groupsize\_zvf\_sweep.\{tsv,json\}}).

\paragraph{S3: GSM8K reward tensors.} Qwen3-8B on the Tinker sampler, 200
GSM8K problems per seed $\times$ seeds $\{42,123,456\}$, $G{=}8$, temperature
1.0, no training (600 groups;
\texttt{experiments/results/tinker\_gsm8k\_zvf\_s*.json}).

\paragraph{S4: out-of-sample ZVF prediction.} Qwen2.5-1.5B-Instruct (vLLM, no
training) on 200 GSM8K training-split prompts: 32 reference samples per prompt,
then three fresh groups per prompt at each $G\in\{4,8,16\}$ with different
sampling seeds
(\texttt{experiments/results/publication\_fixes/zvf\_novelty/};
\texttt{platform\_colab/publication\_fixes/zvf\_analyze.py}).

\paragraph{E1 and E3 pilots.} Qwen2.5-0.5B-Instruct, full-parameter AdamW, on
synthetic two-operand addition (Sections~\ref{sec:p7-e1}
and~\ref{sec:p7-e3}; \texttt{zvf-program/colab-experiments/}).

\paragraph{N2 reward tensors.} Four GRPO-family methods (\texttt{grpo},
\texttt{aero}, \texttt{areal}, \texttt{gift}), one seed, 40 steps $\times$ 16
GSM8K training prompts per step $\times$ $G{=}8$, the same prompts across
methods at each step (\texttt{experiments/results/n2\_reward\_tensor\_resume/};
launcher \texttt{tinker-runs/scripts/n2\_reward\_tensor\_20260704.py}, whose
default checkpoint is Qwen3.5-4B; the tensor files do not record the
checkpoint). N2 is used only for the descriptive count in
Section~\ref{sec:p7-replay-boundary}.

\paragraph{Statistical conventions.} The unit of replication is the seed (or,
for S4, the prompt; intervals there are prompt-level cluster bootstraps with
10{,}000 draws). Training steps and prompt-steps within one trajectory are
never treated as independent replicates. With one or two seeds per arm (E1,
E3, N2) no inferential test is reported. No confirmatory hypothesis family is
declared, so no multiplicity correction is applied, and a non-significant
difference is never read as equivalence.
